\documentclass[10pt,a4paper]{amsart}
\usepackage[margin=19mm]{geometry}
\usepackage{amsmath,amssymb,mathtools}
\usepackage[scr=rsfs,cal=euler]{mathalfa}
\usepackage{xfrac}
\usepackage[unicode,hidelinks,bookmarks=false]{hyperref}
\newcommand{\ie}{i.e.\ }
\newcommand{\cf}{cf.\ }
\newcommand{\resp}{respectively\ }
\newcommand{\Subsection}{\subsection}
\providecommand{\nobreakdash}{\nobreak\hbox{-}}
\setlength{\emergencystretch}{2em}
\multlinegap0pt
\usepackage{amssymb}
\usepackage{enumerate}
\newtheorem{lemma}{Lemma}
\newtheorem{proposition}{Proposition}
\title{Improved bounds on the number of holomorphic maps between compact Riemann surfaces}
\author{Masaharu Tanabe}
\date{Source: arXiv:2605.21137v2 (CC BY 4.0)}
\begin{document}
\maketitle

\noindent\emph{Source and licence.} Improved bounds on the number of holomorphic maps between compact Riemann surfaces. Version arXiv:2605.21137v2 (CC BY 4.0), distributed by arXiv under the Creative Commons Attribution 4.0 International licence, http://creativecommons.org/licenses/by/4.0/, as recorded in the arXiv licence metadata for this exact version and shown on the abstract page https://arxiv.org/abs/2605.21137v2. Full licence notice: \texttt{licenses/arxiv-2605.21137-CC-BY-4.0.txt}.\par\smallskip

\noindent\textit{Editorial scope.} Counted unit: the article's proposition prop (source0.tex lines 535--546). The article's complete original proof of it is delivered with it (source0.tex lines 548--601, one \texttt{proof} environment of 54 lines). No mathematics is written, repaired or re-derived: the statement and the proof are the article's own lines, in the article's own notation, and no retained line is substituted or re-flowed.  Retained uncounted as the source-local context the proof needs: lines 282--290; lines 293--300.  Nothing was omitted from any retained span, so the package carries no omission record.  No retained line contains a citation, so the wrapper carries no bibliography and no bibliography entry.  No retained line contains a figure environment, an image-inclusion command or drawing code, so no image is delivered and the package declares no asset dependency.  Editorial formatting only, in addition: an AMS \texttt{amsart} preamble that restates the article's own declarations and macros used by the retained lines, namely the environments \texttt{lemma}, \texttt{proposition} on separate counters, together with the standard packages those lines need.  The retained environments print in this document as Lemma 1, Proposition 1 in order of appearance, whereas the article prints the same items with the numbering of its own full document.  The complete native source of the article as distributed in the arXiv e-print for this exact version is bundled unchanged as \texttt{sources/201041/source0.tex}.
\begin{lemma}\label{lem1} For an arbitrary $\mathfrak{g}\in End(J(X))$, denoting by $G$ and $^{t}G^{\prime}$ the rational representations of the duals $^{t}\mathfrak{g}$ and $^{t}\mathfrak{g^{\prime}}$, respectively, we have
\begin{equation}
(^{t}Gx,^{t}Gy)_{X}=(^{t}G^{\prime}{^{t}Gx},y)_{X}=(x,^{t}G^{\prime}{^{t}Gy})_{X}, 
\end{equation}
for all $x, y\in\mathbb{R}^{2g}.$ In particular, when $\mathfrak{F}$ is the endomorphism associated with some holomorphic map $f:X\rightarrow Y$ of degree $d$, denoting by $\mathcal{F}$ the rational representation of $\mathfrak{F}$, we have
\begin{equation}
(^{t}\mathcal{F}x,^{t}\mathcal{F}y)_{X}=d(^{t}\mathcal{F}x,y)_{X}. 
\end{equation}
\end{lemma}

 \begin{lemma} \label{lem2} 
 Let $f:X\rightarrow Y$ be a nonconstant holomorphic map of degree $d$, and let $\mathfrak{F}$ be the endomorphism associated with $f$. 
 Then, denoting by $\mathcal{F}$ the rational representation of $\mathfrak{F}$, we have
\begin{equation}
||^{t}\mathcal{F}x||\le d||x||, 
\end{equation}
for an arbitrary vector $x\in\mathbb{R}^{2g}$.
\end{lemma}

\begin{proposition}\label{prop}
Let $f_{i}:X\rightarrow Y_{i}$ be a nonconstant holomorphic map of degree $d_i\leq d$, 
and let $\mathcal{F}_{i}$ be the rational representation of the endomorphism associated with $f_{i}$.  Suppose that, for some $k<2g$,
\begin{equation*}
{}^{t}\mathcal{F}_{i}a_{1}=\cdot\cdot\cdot=^{t}\mathcal{F}_{i}a_{k-1}=0, 
\end{equation*}
and that  $^{t}\mathcal{F}_{i}a_{k}\ne 0$ hold.  Then the number of possible $^{t}\mathcal{F}_{i}a_{k}$
is at most 
$$
(2d+1)^{2g-k+1}.
$$
\end{proposition}

\begin{proof}
In Euclidean space of $n$ dimensions, the volume $V_n(R)$ of the ball of radius $R$ is given by
$$
V_n(R)=\frac{\pi^{\frac{n}{2}}}{\Gamma (\frac{n}{2}+1)}R^n .
$$

Suppose that $^{t}\mathcal{F}_{1}a_{1}\ne 0$ and $^{t}\mathcal{F}_{2}a_{1}\ne 0$.
By Lemma {\ref{lem2}}, 
$$||^{t}\mathcal{F}_i (a_1 )||\le d_i ||a_1 ||=d_i \lambda_1 \quad (i=1,2).
$$
Since $\lambda_1$ is the minimum norm in the lattice, the difference  satisfies
$$||^{t}\mathcal{F}_1 (a_1 )-^{t}\mathcal{F}_2 (a_1 )||\ge \lambda_1
$$ if it is not $0$.
Thus we see that the number of possible $^{t}\mathcal{F}_{i}a_{1}$
is at most 
$$
\frac{V_{2g} \left( d+\frac1{2} \right) }{V_{2g} \left( \frac1{2} \right)}
=(2d+1)^{2g}.
$$

Next, suppose that
$$
{}^{t}\mathcal{F}_{i}a_{1}=0
\qquad\text{and}\qquad
{}^{t}\mathcal{F}_{i}a_{2}\neq 0.
$$
Then
$
{}^{t}\mathcal{F}_{i}a_{2}
$
is orthogonal to \(a_{1}\). Indeed, by Lemma~\ref{lem1},
$$
d\,({}^{t}\mathcal{F}_{i}a_{2},a_{1})_X
=
({}^{t}\mathcal{F}_{i}a_{2},{}^{t}\mathcal{F}_{i}a_{1})_X
=
0.
$$
Hence
$
{}^{t}\mathcal{F}_{i}a_{2}
$
lies in the \((2g-1)\)-dimensional subspace orthogonal to $a_{1}$.
Since \(\lambda_{2}\) is the first successive minimum of the induced
sublattice in this subspace, 
the number of possible $^{t}\mathcal{F}_{i}a_{2}$
is at most 
$$
\frac{V_{2g-1} \left( d+\frac1{2} \right) }{V_{2g-1} \left( \frac1{2} \right)}
=(2d+1)^{2g-1}.
$$

Proceeding by induction, we obtain the conclusion.
\end{proof}

\end{document}
